% AA-51 — klasyfikacja typów izomorfizmu, nie klasyfikacja sprzężoności.
\chapter{Skończone podgrupy grupy przekształceń Möbiusa}
\chaptermark{Skończone podgrupy PGL(2)}
\label{aa:ch:finite-pgl2}

W rozdziale~\ref{aa:ch:kovacic-finite-groups} wykorzystaliśmy
klasyfikację skończonych prymitywnych grup rzutowych. Teraz
udowodnimy klasyfikację typów izomorfizmu wszystkich skończonych
podgrup $\operatorname{PGL}_2(\mathbb C)$ oraz wyjaśnimy
pochodzenie liczb $4,6,12$. Nie dowodzimy tutaj jednoznaczności
zanurzeń z dokładnością do sprzężenia ani pełnego twierdzenia
o listach lokalnych algorytmu Kovacica.

\section{Działanie na prostej rzutowej}

\begin{definicja}[Grupa rzutowa]
\label{aa:def:pgl2-projective}
Niech $\mathbb P^1(\mathbb C)$ będzie zbiorem prostych
wektorowych w $\mathbb C^2$. Definiujemy
$\operatorname{PGL}_2(\mathbb C)=
\operatorname{GL}_2(\mathbb C)/\{\lambda I:\lambda\in\mathbb C^\times\}$.
Klasa macierzy $A$ działa przez $[v]\mapsto[Av]$.
Współrzędna $z=v_1/v_2$ identyfikuje prostą rzutową
z $\mathbb C\cup\{\infty\}$, a działanie macierzy
$\left(\begin{smallmatrix}a&b\\c&d\end{smallmatrix}\right)$
z przekształceniem $z\mapsto(az+b)/(cz+d)$,
z naturalnym przedłużeniem w biegunie i w nieskończoności.
\end{definicja}

Działanie jest wierne: macierz zachowująca wszystkie proste
zachowuje proste rozpięte przez $e_1,e_2,e_1+e_2$,
więc jest skalarna. Ponadto każda macierz odwracalna
może być pomnożona przez skalar tak, aby jej wyznacznik
wynosił $1$. Stąd $\operatorname{PGL}_2(\mathbb C)
\simeq\operatorname{SL}_2(\mathbb C)/\{\pm I\}
=\operatorname{PSL}_2(\mathbb C)$.

\begin{lemat}[Punkty stałe i stabilizatory]
\label{aa:lem:pgl2-fixed}
Niech $G$ będzie skończoną podgrupą
$\operatorname{PGL}_2(\mathbb C)$. Każdy $g\in G\setminus\{1\}$
ma dokładnie dwa punkty stałe w $\mathbb P^1(\mathbb C)$,
a stabilizator $G_p$ każdego punktu $p$ jest cykliczny.
\end{lemat}
\begin{proof}
Jeżeli $g$ ma rząd $m>1$ i reprezentant $A$,
to $A^m=\lambda I$ dla pewnego $\lambda\ne0$.
Wielomian $T^m-\lambda$ ma różne pierwiastki,
więc $A$ jest diagonalizowalna. Gdyby miała tylko
jedną wartość własną, byłaby skalarna. Zatem ma dokładnie
dwie proste własne, będące punktami stałymi $g$.

Po zmianie współrzędnych możemy przyjąć $p=\infty$.
Wówczas każdy $h\in G_p$ działa przez $z\mapsto a_hz+b_h$.
Liczba $c=|G_p|^{-1}\sum_{h\in G_p}h(0)$ jest wspólnym
punktem stałym: dla $k\in G_p$ afiniczność daje
$k(c)=|G_p|^{-1}\sum_h(kh)(0)=c$.
Po przesunięciu $c$ do zera grupa działa przez mnożenia
$z\mapsto a_hz$ i zanurza się w $\mathbb C^\times$.
Jeśli jej rząd wynosi $s$, każdy $a_h$ jest pierwiastkiem
równania $T^s=1$. Ponieważ istnieje dokładnie $s$ takich
pierwiastków i tworzą grupę cykliczną, stabilizator
jest cykliczny.
\end{proof}

W szczególności element ustalający trzy różne punkty
jest identycznością. To proste spostrzeżenie będzie
również użyte przy konstrukcji wiernych reprezentacji
permutacyjnych.

\section{Równanie orbit wyjątkowych}

\begin{twierdzenie}[Możliwe dane stabilizatorów]
\label{aa:tw:pgl2-signatures}
Niech $G\le\operatorname{PGL}_2(\mathbb C)$ będzie
skończona, $N=|G|>1$. Zbiór punktów o nietrywialnym
stabilizatorze rozkłada się na $r$ orbit $O_1,\ldots,O_r$.
Wybierzmy $p_i\in O_i$ i oznaczmy $m_i=|G_{p_i}|$.
Po zmianie numeracji przyjmujemy $2\le m_1\le\cdots\le m_r$.
Zachodzi
\begin{equation}\label{aa:eq:pgl2-orbit-count}
 \sum_{i=1}^r\left(1-\frac1{m_i}\right)=2-\frac2N .
\end{equation}
Możliwe są dokładnie następujące dane liczbowe:
\[
\begin{array}{c|c}
(m_1,\ldots,m_r)&N\\ \hline
(N,N)&N\\
(2,2,n),\ n\ge2&2n\\
(2,3,3)&12\\
(2,3,4)&24\\
(2,3,5)&60
\end{array}
\]
\end{twierdzenie}
\begin{proof}
Zbiór punktów o nietrywialnym stabilizatorze jest
sumą zbiorów punktów stałych elementów $g\ne1$.
Na mocy Lematu~\ref{aa:lem:pgl2-fixed} jest skończony;
jest niezmienniczy względem $G$, ponieważ
$G_{hp}=hG_ph^{-1}$ dla $h\in G$.
Zliczamy pary $(g,p)$, gdzie $g\ne1$ i $g(p)=p$.
Z Lematu~\ref{aa:lem:pgl2-fixed} ich liczba to $2(N-1)$.
Orbita o stabilizatorze rzędu $m_i$ zawiera $N/m_i$
punktów i każdy z nich jest ustalany przez $m_i-1$
nietrywialnych elementów. Podzielenie równości liczby
par przez $N$ daje~\eqref{aa:eq:pgl2-orbit-count}.
Każdy składnik lewej strony jest co najmniej $1/2$,
więc $r\le3$. Prawa strona jest co najmniej $1$,
a pojedynczy składnik jest mniejszy od $1$, zatem $r\ge2$.
Dla $r=2$ otrzymujemy $1/m_1+1/m_2=2/N$.
Ponieważ $m_i\le N$, obie liczby muszą równać się $N$.

Dla $r=3$ mamy $1/m_1+1/m_2+1/m_3=1+2/N>1$.
Stąd $m_1=2$; inaczej suma byłaby nie większa od $1$.
Jeśli $m_2=2$, równanie daje $N=2m_3$.
Jeżeli $m_2\ge4$, suma znowu jest nie większa od $1$.
Pozostaje $m_2=3$, a wtedy $1/m_3>1/6$,
czyli $m_3=3,4,5$. Podstawienie daje odpowiednio
$N=12,24,60$. Istnienie grup realizujących te dane
wykażemy poniżej.
\end{proof}

\section{Identyfikacja grup}

\begin{lemat}[Podgrupy indeksu dwa grup symetrycznych]
\label{aa:lem:pgl2-index2}
Dla $k\ge2$ jedyną podgrupą indeksu $2$ grupy $S_k$
jest grupa alternująca $A_k$.
\end{lemat}
\begin{proof}
Podgrupa indeksu $2$ jest normalna i jest jądrem
surjektywnego homomorfizmu do grupy $C_2$.
Wszystkie transpozycje są sprzężone, więc mają ten sam
obraz. Generują $S_k$, a homomorfizm jest surjektywny;
obraz transpozycji jest więc nietrywialny. Homomorfizm
pokrywa się z homomorfizmem znaku, którego jądro to $A_k$.
\end{proof}

\begin{twierdzenie}[Klasyfikacja typów izomorfizmu]
\label{aa:tw:pgl2-classification}
Każda skończona podgrupa $G$ grupy
$\operatorname{PGL}_2(\mathbb C)$ jest izomorficzna
z jedną z grup: cykliczną $C_n$ ($n\ge1$),
dwuścienną $D_n$ rzędu $2n$ ($n\ge2$),
$A_4$, $S_4$ lub $A_5$.
W przypadku $n=2$ grupa $D_2$ jest grupą Kleina.
\end{twierdzenie}
\begin{proof}
Grupa trywialna jest $C_1$. Dla $N>1$ korzystamy
z Twierdzenia~\ref{aa:tw:pgl2-signatures}.

\emph{Dane $(N,N)$.} Grupa ustala punkt, więc sama
jest jego cyklicznym stabilizatorem.

\emph{Dane $(2,2,n)$.} Orbita odpowiadająca stabilizatorowi
rzędu $n$ ma dwa punkty. Po zmianie współrzędnych
są to $0,\infty$. Jądro działania na tej parze jest
stabilizatorem jednego punktu, cyklicznym rzędu $n$,
generowanym przez $a:z\mapsto\zeta z$, gdzie $\zeta$
jest pierwotnym pierwiastkiem $n$-tego stopnia z jedności.
Element $b$ zamieniający punkty ma postać $z\mapsto c/z$,
$c\ne0$. Stąd $b^2=1$ i $bab^{-1}=a^{-1}$.
Wszystkie elementy są postaci $a^j$ albo $a^jb$,
$0\le j<n$, co daje grupę dwuścienną rzędu $2n$.

\emph{Dane $(2,3,3)$.} Jedna z orbit o stabilizatorze
rzędu $3$ ma cztery punkty. Działanie na niej jest wierne,
bo element ustalający cztery punkty jest identycznością.
Obraz w $S_4$ ma rząd $12$, a więc jest $A_4$
na mocy Lematu~\ref{aa:lem:pgl2-index2}.

\emph{Dane $(2,3,4)$.} Istnieje osiem punktów
o stabilizatorze rzędu $3$. Każda podgrupa rzędu $3$
ustala dokładnie dwa punkty (są to punkty stałe jej generatora),
a ich stabilizatory są właśnie tą podgrupą: w pozostałych
orbitach rzędy stabilizatorów nie są podzielne przez $3$.
Zbiory tych dwóch punktów dla różnych podgrup są rozłączne,
bo każdy punkt z tej orbity wyznacza swój stabilizator.
Otrzymujemy cztery podgrupy Sylowa rzędu $3$.
Działanie przez sprzężenie na tych czterech podgrupach
jest przechodnie. Jego jądro $K$ zawiera się w normalizatorze
jednej z nich, rzędu $24/4=6$.
Gdyby $3$ dzieliło $|K|$, twierdzenie Cauchy'ego
i normalność $K$ wymusiłyby zawieranie wszystkich czterech
podgrup rzędu $3$, a więc co najmniej dziewięciu elementów.
To niemożliwe. Zatem $|K|=1$ lub $2$.
W drugim przypadku jedyny nietrywialny element $h\in K$
byłby centralny. Dla dowolnego elementu $a$ rzędu $3$,
element $h$ zachowuje dwa punkty stałe $a$.
Nie może ich zamieniać: po przeniesieniu do $0,\infty$
zamiana sprzęga mnożenie $a$ z jego odwrotnością,
a $a\ne a^{-1}$. Stąd $h$ ustala oba punkty stałe
każdej podgrupy rzędu $3$, czyli osiem różnych punktów,
wbrew Lematowi~\ref{aa:lem:pgl2-fixed}.
Zatem $K=1$ i obraz w $S_4$ ma rząd $24$, czyli jest $S_4$.

\emph{Dane $(2,3,5)$.} Trzy wyjątkowe orbity mają
odpowiednio $30,20,12$ punktów.
Każdy nietrywialny element należy do stabilizatora
jednego ze swoich punktów stałych, więc jego rząd
to $2,3$ albo $5$.
Każda inwolucja ustala dwa punkty z pierwszej orbity,
a każdy z jej punktów wyznacza jedną inwolucję.
Jest zatem $15$ inwolucji; są sprzężone, ponieważ
punkty tworzą jedną orbitę. Centralizator inwolucji
ma rząd $60/15=4$.
Podgrupa Sylowa rzędu $4$ nie ma elementu rzędu $4$,
więc każdy nietrywialny element ma rząd $2$.
Dla dowolnych jej elementów $a,b$ mamy
$ab=(ab)^{-1}=b^{-1}a^{-1}=ba$.
Jeśli $a,b$ są różnymi nietrywialnymi elementami,
to $1,a,b,ab$ są parami różne i podgrupa jest
izomorficzna z $C_2\times C_2$.
Każda inwolucja należy do takiej podgrupy, a każda
taka podgrupa zawierająca $h$ zawiera się w $C_G(h)$,
czyli jest równa $C_G(h)$. Inwolucje są więc rozdzielone
na rozłączne trójki i istnieje pięć podgrup Sylowa rzędu $4$.
Działanie przez sprzężenie na nich jest przechodnie;
jego jądro $K$ ma rząd dzielący $60/5=12$.
Gdyby $2$ dzieliło $|K|$, normalność i twierdzenie
Cauchy'ego wymusiłyby obecność wszystkich $15$ inwolucji,
co jest niemożliwe. Jeśli $3$ dzieli $|K|$, to $K$
zawiera wszystkie podgrupy Sylowa rzędu $3$.
Jest ich dziesięć, gdyż dwadzieścia punktów drugiej
orbity odpowiada parami ich punktom stałym.
Daje to co najmniej $21$ elementów $K$, również niemożliwe.
Zatem $K=1$. Wierny obraz grupy rzędu $60$ w $S_5$
ma indeks $2$, więc jest $A_5$.
\end{proof}

\section{Realizacja wszystkich typów}

\begin{przyklad}[Grupy cykliczne i dwuścienne]
\label{aa:prz:pgl2-cyclic-dihedral}
Mnożenia $z\mapsto\zeta^jz$ dla $\zeta=\exp(2\pi i/n)$
realizują $C_n$. Dołączenie $z\mapsto1/z$ dla $n\ge2$
daje dokładnie $2n$ różnych przekształceń, o relacjach
z dowodu Twierdzenia~\ref{aa:tw:pgl2-classification}.
Dla $n=2$ są to $z,-z,1/z,-1/z$.
\end{przyklad}

\begin{lemat}[Obroty jako przekształcenia rzutowe]
\label{aa:lem:pgl2-rotations}
Grupa $\operatorname{SO}(3)$ jest izomorficzna z podgrupą
$\operatorname{PGL}_2(\mathbb C)$ będącą obrazem
$\operatorname{SU}(2)$.
\end{lemat}
\begin{proof}
Niech $V$ będzie rzeczywistą przestrzenią hermitowskich
macierzy $2\times2$ o śladzie zero.
Macierz
$\left(\begin{smallmatrix}z&x-iy\\x+iy&-z\end{smallmatrix}\right)$
identyfikuje $V$ z $\mathbb R^3$;
jej ujemny wyznacznik to $x^2+y^2+z^2$.
Sprzężenie $H\mapsto UHU^{-1}$ przez $U\in\operatorname{SU}(2)$
zachowuje tę normę i daje homomorfizm do $\operatorname O(3)$.
Macierze $\operatorname{SU}(2)$ mają postać
$\left(\begin{smallmatrix}a&b\\-\overline b&\overline a\end{smallmatrix}\right)$,
$|a|^2+|b|^2=1$. Tworzą więc przestrzeń homeomorficzną
ze sferą $S^3$, która jest spójna.
Wyznacznik indukowanego odwzorowania jest ciągły,
przyjmuje tylko wartości $\pm1$ i w identyczności wynosi $1$,
więc obraz leży w $\operatorname{SO}(3)$.
Jądro składa się z $\pm I$: macierz komutująca z
$\operatorname{diag}(1,-1)$ jest diagonalna, a komutowanie
z $\left(\begin{smallmatrix}0&1\\1&0\end{smallmatrix}\right)$
wymusza równość wpisów diagonalnych; wyznacznik to $1$.

Sprzężenie przez $\operatorname{diag}(e^{-it/2},e^{it/2})$
pozostawia $z$ bez zmiany, a współrzędną $x+iy$
mnoży przez $e^{it}$, daje więc obrót o kąt $t$
wokół osi $z$.
Każda jednostkowa macierz $H\in V$ ma wartości własne
$1,-1$; ortonormalną bazę własną można wybrać tak,
aby macierz zmiany bazy miała wyznacznik $1$ (wystarczy
pomnożyć jeden wektor przez skalar modułu $1$).
Stąd każda oś jednostkowa jest obrazem osi $z$
przez element obrazu $\operatorname{SU}(2)$.
Każdy obrót w $\operatorname{SO}(3)$ jest obrotem wokół osi:
rzeczywista macierz ortogonalna wymiaru trzy o wyznaczniku $1$
ma wartość własną $1$, a na prostopadłej płaszczyźnie
działa przez dwuwymiarowy obrót.
Sprzęgając obrót wokół osi $z$ otrzymujemy więc każdy
element $\operatorname{SO}(3)$. Ostatecznie
$\operatorname{SO}(3)\simeq\operatorname{SU}(2)/\{\pm I\}$,
a ten iloraz zanurza się w $\operatorname{PGL}_2(\mathbb C)$.
\end{proof}

\begin{przyklad}[Czworościan i sześcian]
\label{aa:prz:pgl2-tetra-cube}
Na zbiorze czterech wierzchołków
$(1,1,1),(1,-1,-1),(-1,1,-1),(-1,-1,1)$
działają obroty otrzymane z cyklicznych permutacji
trzech współrzędnych i zmian znaków o iloczynie $1$.
Te $3\cdot4=12$ macierzy tworzy grupę.
Działanie jest wierne, ponieważ wierzchołki rozpinają
$\mathbb R^3$, więc obraz ma indeks dwa w $S_4$
i jest $A_4$.

Macierze permutujące osie współrzędnych i zmieniające znaki,
o wyznaczniku $1$, tworzą grupę $6\cdot4=24$ obrotów
sześcianu. Działają na czterech prostych zawierających
przekątne sześcianu. Element jądra zachowuje proste
o wektorach kierunkowych z poprzedniego akapitu.
Dowolne trzy z nich są liniowo niezależne, a ich suma
jest przeciwna czwartemu. Wektory są zatem wektorami
własnymi o tej samej wartości własnej $\pm1$.
Wyznacznik $1$ wymusza identyczność.
Działanie jest wierne i grupa jest $S_4$.
\end{przyklad}

\begin{przyklad}[Dwudziestościan]
\label{aa:prz:pgl2-icosahedron}
Niech $\phi=(1+\sqrt5)/2$ i niech $X$ składa się
z dwunastu wektorów
$(0,\pm1,\pm\phi),(\pm1,\pm\phi,0),(\pm\phi,0,\pm1)$.
Grupa $G_X=\{R\in\operatorname{SO}(3):R(X)=X\}$
jest skończona, bo $X$ rozpina $\mathbb R^3$ i działanie
na $X$ jest wierne. Jest przechodnia na $X$:
cykliczne permutacje współrzędnych oraz zmiany znaków
o iloczynie $1$ pozwalają przejść między dowolnymi
wierzchołkami (znak zerowej współrzędnej można wybrać
tak, aby zachować wyznacznik).

Dla $v=(0,1,\phi)$ pięć wektorów o iloczynie skalarnym
z $v$ równym $\phi$ to, w kolejności cyklicznej,
\[
 (0,-1,\phi),\quad(\phi,0,1),\quad(1,\phi,0),
 \quad(-1,\phi,0),\quad(-\phi,0,1).
\]
Pozostałe wierzchołki to $v,-v$ i przeciwne do tych
pięciu wektorów. Oznaczmy te pięć wektorów przez $w_0,\ldots,w_4$,
z indeksami liczonymi modulo $5$. Wszystkie mają
kwadrat normy $1+\phi^2$, a z $\phi^2=\phi+1$ wynika
\[
 \langle w_i,w_j\rangle=
 \begin{cases}
  1+\phi^2,& i=j,\\
  \phi,& i-j\equiv\pm1\pmod 5,\\
  -\phi,& i-j\equiv\pm2\pmod 5.
 \end{cases}
\]
Możemy więc skonstruować obrót bez odwoływania się
do wizualnej regularności pięciokąta.
Niech $B=(v,w_0,w_1)$ oraz $C=(v,w_1,w_2)$ będą
macierzami o wskazanych kolumnach. Mamy
$\det B=2\phi^2\ne0$ oraz $B^{\mathsf T}B=C^{\mathsf T}C$.
W szczególności $C$ jest odwracalna. Odwzorowanie
$R=CB^{-1}$ jest ortogonalne, bo
\[
 R^{\mathsf T}R
 =B^{-\mathsf T}C^{\mathsf T}CB^{-1}=I .
\]
Z definicji $Rv=v$, $Rw_0=w_1$, $Rw_1=w_2$.
Dla każdego $i$ wektory $Rw_i$ i $w_{i+1}$ mają
te same iloczyny skalarne z $v,w_1,w_2$:
dla pierwszego wektora wynika to z ortogonalności $R$,
a dla drugiego z podanej tabeli iloczynów skalarnych.
Ponieważ $v,w_1,w_2$ tworzą bazę, wynika stąd
$Rw_i=w_{i+1}$. Zatem $R^5=I$, gdyż $R^5$ ustala
wektory bazy $B$, oraz $R\ne I$.
Rząd $R$ wynosi $5$, więc $(\det R)^5=1$ i $\det R=1$.
W szczególności $R\in G_X$; zachowuje także przeciwne
wektory oraz $\pm v$. Stabilizator $v$ ma zatem
co najmniej pięć elementów. Każdy $Q\in (G_X)_v$
zachowuje iloczyn skalarny z $v$, zatem $Qw_0$
należy do zbioru $\{w_0,\ldots,w_4\}$.
Jeśli $Q_1w_0=Q_2w_0$, to obrót $Q_2^{-1}Q_1$
ustala dwa liniowo niezależne wektory $v,w_0$.
Ustala więc punktowo rozpiętą przez nie płaszczyznę,
a ponieważ jest ortogonalny i ma wyznacznik $1$,
ustala również jej prostopadłą prostą.
Stąd $Q_1=Q_2$. Stabilizator ma więc najwyżej
pięć elementów.
Stąd $|G_X|=12\cdot5=60$.

Przechodniość daje sześć różnych osi przez pary
$\{w,-w\}$, $w\in X$, każda z podgrupą obrotów rzędu $5$.
Są to różne podgrupy, bo nietrywialny obrót ma jedną oś.
Grupa cykliczna rzędu $60$ i grupa dwuścienna rzędu $60$
mają tylko jedną podgrupę rzędu $5$ (w dwuściennej
wszystkie elementy nieparzystego rzędu leżą w grupie
obrotów). Z Twierdzenia~\ref{aa:tw:pgl2-classification}
i Lematu~\ref{aa:lem:pgl2-rotations} wynika $G_X\simeq A_5$.
\end{przyklad}

\begin{uwaga}[Realizowalność danych stabilizatorów]
\label{aa:uw:pgl2-signatures-realized}
Powyższe konstrukcje realizują również każdą sygnaturę
z Twierdzenia~\ref{aa:tw:pgl2-signatures}.
Nietrywialna grupa cykliczna realizuje $(N,N)$,
a grupa dwuścienna rzędu $2n$, $n\ge2$, realizuje
$(2,2,n)$: zachowuje parę $\{0,\infty\}$ i nie ma
wspólnego punktu stałego. Dla $n=2$ wszystkie trzy
wyjątkowe orbity mają dwa punkty.
W grupach $A_4,S_4,A_5$ możliwe rzędy elementów
to odpowiednio $1,2,3$; $1,2,3,4$ oraz $1,2,3,5$
(wynika to z rozkładów permutacji na rozłączne cykle).
W szczególności grupy te nie są cykliczne ani dwuścienne:
grupy cykliczne tych samych rzędów zawierałyby element
rzędu $12,24,60$, a grupy dwuścienne element rzędu
$6,12,30$. Zatem jedynymi możliwymi danymi stabilizatorów
są odpowiednio $(2,3,3)$, $(2,3,4)$ i $(2,3,5)$.
\end{uwaga}

\section{Powrót do różniczkowej teorii Galois}

\begin{wniosek}[Prymitywny przypadek i długości orbit]
\label{aa:wn:pgl2-primitive-orbits}
Niech $H\le\operatorname{SL}_2(\mathbb C)$ będzie skończona,
a jej naturalna reprezentacja na $\mathbb C^2$ prymitywna
w sensie Definicji~\ref{aa:def:kov3-primitive}.
Obraz rzutowy $G$ grupy $H$ jest izomorficzny z $A_4,S_4$
lub $A_5$, a najmniejsze długości orbit na
$\mathbb P^1(\mathbb C)$ wynoszą odpowiednio $4,6,12$.
\end{wniosek}
\begin{proof}
Gdy obraz rzutowy jest trywialny lub cykliczny,
grupa $H$ zachowuje prostą (dla nietrywialnej grupy
cyklicznej wybieramy punkt stały generatora).
Gdy obraz jest dwuścienny, zachowuje dwuelementową
orbitę, czyli nieuporządkowaną parę różnych prostych,
których suma jest całą przestrzenią. Oba przypadki
przeczą prymitywności.
W pozostałych trzech przypadkach minimalne długości
orbit wynoszą $N/\max_i m_i$, bo orbita zwykła ma
$N$ punktów. Z danych $(2,3,3)$, $(2,3,4)$, $(2,3,5)$
otrzymujemy $12/3=4$, $24/4=6$, $60/5=12$.
\end{proof}

Liczby te wyjaśniają geometryczne pochodzenie parametrów
z rozdziału~\ref{aa:ch:kovacic-finite-groups}.
Sama klasyfikacja grup i obliczenie długości orbit
nie dowodzą jednak lokalnych warunków dla potencjału
ani rekurencji koniecznej w algorytmie Kovacica.
